\section{Einleitung}
Die modernen Stromnetze werden durch die stärkere Digitalisierung und die Menge an dezentralen Energiequellen immer komplexer. Sie entwickeln sich zu sogenannten Smart Grids. Um einen effizienten und stabilen Netzbetrieb zu garantieren, kombinieren sie dabei die physische Strominfrastruktur mit Informations- und Kommunikationstechnologien (Gungor et al., 2013). Denn schwankende Erzeugung und dynamische Lastsituationen können das Netz, z.B. durch die stetig wachsende Anzahl an einspeisenden PVC-Anlagen, überlasten. Gleichzeitig erhöht diese enge Vernetzung jedoch auch die Verwundbarkeit des Systems. Ausfälle einzelner Komponenten, Überlastsituationen, oder gezielte Angriffe können sich aufgrund durch wechselseitige Abhängigkeiten schnell ausbreiten und damit kritische Stromversorgungsstörungen verursachen. Tatsächlich zeigen frühere großflächige Netzausfälle bereits deutliche Beispiele solcher Gefahren (Buldyrev et al. 2010).
Vor diesem Hintergrund rückt das Konzept der Resilienz verstärkt in den Fokus. Resilienz definiert, wie gut ein System unter Belastung weiterarbeitet oder sich danach wieder erholt. Damit gilt es als entscheidender Maßstab für künftige Energiesysteme (Holling, 1973).
Zur Erforschung der Strukturmerkmale von Smart Grids bei Störungen bietet die Theorie komplexer Netze eine nützliche Grundlage. Sie ermöglicht es, Energienetze als graphbasierte Modelle darzustellen. Es können kritische Knoten und Verbindungen identifiziert werden und die Robustheit und Resilienz können gegenüber zufälligen Ausfällen und gezielten Angriffen quantitativ bewertet werden (Albert, Jeong & Barabási, 2000). Unterschiedliche Netzwerktopologien, etwa zufällige Netzwerke nach Erdős und Rényi, Small-World-Strukturen nach Watts und Strogatz oder skalenfreie Netzwerke nach Barabási und Albert zeigen charakteristische Resilienz- und Verwundbarkeitsmuster auf. Diese Typen unterscheiden sich klar im Grad ihrer Widerstandskraft sowie Schwachstellenverteilung, was direkt für den Entwurf echter Energieversorgungsnetze genutzt werden kann (Newman, 2018). Vor diesem Hintergrund vielfältiger Analysemöglichkeiten entsteht daher folgende Fragestellungen: Welche konkreten Merkmale beeinflussen tatsächlich die Stabilität dieser Netze? Wie lässt sich diese durch gezielte Maßnahmen steigern?
Zweck dieser Studie ist die systematische Untersuchung der Resilienz von Smart Grids anhand netzwerkbasierter Verfahren und algorithmischer Auswertungen. Dafür werden Stromnetze als komplexe Netze dargestellt, unterschiedliche Störverläufe nachgebildet sowie Maßzahlen für Widerstandsfähigkeit ermittelt. Mithilfe dieser Ergebnisse lassen sich dann konkrete Gestaltungselemente erkennen, auch Hinweise auf mögliche Verbesserungen bei der Stabilität sollten sich erkennen lassen. Damit soll ein Teil zum besseren Erfassen der Netzrobustheit geleistet werden und gleichzeitig Impulse für den Aufbau flexibler, belastbarer und langfristig tragfähiger Systeme geliefert werden.
 